\chapter{Molecular biology for computing readers}
\label{book:appendixA}

This appendix collects the biology that the genomic chapters assume. It explains what
the symbols in a sequence file refer to, why orientation and coordinates are part of
the data contract, and which experimental measurements supply the labels that models
predict. It is a reader's guide for computing specialists, not a laboratory manual. It
contains no procedures, and it makes no clinical interpretation of variants. The
standard textbook treatment is Alberts and colleagues~\cite{albertscontrol}.

\section{The molecule behind the string}
DNA is a polymer of nucleotides. Each nucleotide joins a deoxyribose sugar, a phosphate
group and one of four bases: adenine (A), cytosine (C), guanine (G) or thymine (T).
Phosphodiester linkages join the $3'$ carbon of one sugar, through a phosphate group, to
the $5'$ carbon of the next.
A single strand therefore has a chemical direction, and a sequence is written by
convention from its $5'$ end to its $3'$ end.

In the double helix two strands pair \emph{antiparallel}: one runs $5'\to3'$ left to
right, the other $3'\to5'$. A pairs with T and G pairs with C. The second strand is
therefore determined by the first, and reading it in its own $5'\to3'$ direction gives
the \emph{reverse complement} $R(x)$. Chapters~5, 18, 21 and~24 build their strand
symmetries on exactly this fact. Complementing without reversing describes the partner
strand read in the wrong direction. It is a common bug, not a biological operation.

\BackPlate{duplex-dogma}{Panel A: a short duplex. Each strand carries a half-arrowhead
at its $3'$ end; grey rungs are base pairs. Read $5'\to3'$, the lower strand is the
reverse complement of the upper one. Panel B: the flow from DNA to RNA to protein. RNA
polymerase reads the template strand $3'\to5'$ and synthesizes RNA $5'\to3'$, so the RNA
has the sequence of the coding strand with U in place of T. The ribosome reads codons
$5'\to3'$. Regulatory factors (Chapter~11) act at several of these steps. Shapes are
schematic.}{book:duplex}

\section{Genomes, assemblies and coordinates}
A \emph{genome} is the complete DNA of an organism, packaged into chromosomes. A
\emph{reference assembly} is a versioned consensus sequence used as a coordinate system,
and positions are meaningful only relative to a named assembly. The same variant can
have different integer coordinates in two assemblies. Converting between assemblies is
a separately versioned mapping that can fail or be ambiguous (Chapter~48). A
\emph{locus} is a region in that coordinate system. It may contain a gene, a regulatory
element or neither.

File formats disagree about how positions are counted, which is a classic source of
off-by-one errors:
\begin{center}\small
\begin{tabular}{@{}P{0.22\linewidth}P{0.34\linewidth}P{0.36\linewidth}@{}}
\toprule
Format & Position convention & Example\\\midrule
BED & zero-based start, end excluded~\cite{ucscformats} & the first 100 bases are \texttt{0 100}\\
VCF & one-based \texttt{POS}~\cite{vcfspec} & the first base has \texttt{POS}~1\\
GFF3 & one-based start and end, start $\le$ end~\cite{gff3spec} & the first base starts at 1\\
This book & zero-based, half-open $[a,b)$ & the first 100 bases are $[0,100)$\\
\bottomrule
\end{tabular}
\end{center}
Annotations also carry a strand (\texttt{+} or \texttt{-}). A feature on the minus
strand is read from the reverse complement. Its coordinates map as $[a,b)\mapsto
[L-b,L-a)$ inside a length-$L$ window (Chapters~21 and~24). The symbol \texttt{N}
marks an unknown base. It is an observation symbol, not a fifth nucleotide, and the
book's reference code rejects or masks it explicitly rather than mapping it to A.

\section{From gene to function}
Transcription copies a stretch of DNA into RNA. RNA polymerase reads the template
strand and synthesizes RNA $5'\to3'$. In eukaryotes the initial transcript is
processed: introns are removed by splicing, and the mature messenger RNA (mRNA) leaves
the nucleus. Translation reads mRNA in three-base codons and assembles a protein. Many
RNAs are never translated and act directly. ``Expression'' can therefore mean RNA
abundance, protein abundance or a downstream activity. Chapter~11 shows why these must
be named separately, and why protein data alone may be unable to identify the rates
that produced them.

Expression is regulated at many steps~\cite{albertscontrol}. \emph{Promoters} are
regions where the transcription machinery assembles. \emph{Enhancers} and other
regulatory elements, sometimes far from the gene, bind \emph{transcription factors}
that recognize short sequence motifs. Whether a region is accessible depends on
\emph{chromatin}, the packaging of DNA with histone proteins and its chemical
modifications. In three dimensions, distant elements can come into contact. These
facts motivate the book's interest in long context, distal access and stateful
regulation. Remember throughout that a model's sensitivity to a motif is not proof that
a protein binds it in a cell (Chapters~17 and~48).

\section{Where the labels come from}
Genomic models are trained and judged against measurements, and each kind of
measurement defines a different prediction task. In outline:
\begin{description}
\setlength{\itemsep}{1pt}
\item[RNA sequencing] counts transcripts and supplies expression levels.
\item[Chromatin accessibility assays] mark regions open to cutting enzymes.
\item[Immunoprecipitation-based assays] report where a chosen protein or histone mark
associates with DNA.
\item[Chromosome conformation assays] report which regions are near each other in
three dimensions.
\item[Massively parallel reporter assays] measure the regulatory activity of many
designed sequences in one experiment.
\end{description}
Every assay has a cell type, condition, resolution, noise model and processing pipeline,
and a label is only as meaningful as that provenance (Chapters~2, 25 and~48). A
model that predicts one assay in one cell type has not thereby learned regulation in
general.

\section{Variation and relatedness}
Genomes differ. A single-nucleotide variant (SNV) replaces one base. An insertion or
deletion (indel) changes length, and structural variants rearrange larger segments.
A variant is described by assembly, position, reference allele and alternate allele.
Under reverse complementation both alleles are complemented and the position is
mirrored (Chapter~48). Sequences also share ancestry. Homologous regions within and
between species can be nearly identical, so a random split of genomic windows can put
near-copies on both sides of a train/test boundary. This is why the book closes
leakage relations transitively before splitting (Chapters~2, 25 and~48).

\section{What this appendix does not cover}
Laboratory methods, organism engineering and clinical genetics are outside the book.
Where a chapter proposes biological validation, it names the hypothesis, intervention,
controls and replication unit that specialists would need. It does not describe how to
perform the experiment.
